\documentclass[../main.tex]{subfiles}

\begin{document}

\section{Context of the activity}

An earthquake is the sudden release of energy accumulated within the Earth's tectonic plates. 
This generates vibrations that propagate in the form of seismic waves. 
Seismologists use scales such as the Richter scale to measure the strength of these phenomena.
Table~\ref{tab:sismos-mexico} shows the dates, origin and magnitude of some earthquakes that have happened in Mexico.

\begin{table}[ht!]
\centering
\begin{tabular}{c c c}
\toprule
\textbf{Date} & \textbf{Place} & \textbf{Magnitude} \\
\midrule
28 mar 1787 & Oaxaca coastline & $8.6$ \\
19 sep 1985 & Michoacán / CDMX & $8.1$ \\
7 sep 2017 & Gulf of Tehuantepec & $8.2$ \\
19 sep 2017 & Morelos-Puebla & 7.1 \\
6 jul 1964 & Tierra Caliente & $7.2-7.3$ \\
23 ago 1965 & Oaxaca & $7.3-7.5$ \\
20 mar 2012 & South of Ometepec, Guerrero & $7.5-7.7$ \\
\bottomrule
\end{tabular}
    \caption{Major earthquakes in Mexico and the energy they release}\label{tab:sismos-mexico}
\end{table}

The Richter scale is based on a logarithmic function that relates the amplitude of the waves to a numerical magnitude. 
For instance, an earthquake measuring 4.0 on the Richter scale may be perceived as a slight vibration indoors, whereas a magnitude of 6.0 could result in moderate damage to vulnerable structures. 
In contrast, an earthquake measuring 8.0 or higher, such as the one that struck Michoacán in 1985, can be devastating.

Using mathematical models such as the Gutenberg-Richter equation enables us to calculate the energy released during an earthquake and make quantitative comparisons between events. 
According to Gutenberg-Richter's law, which describes the seismic distribution of magnitude and frequency in Mexico, the typical constants adjusted for the central region of the country\footnote{Equation retrieve from the journal of the Mexican geophysical union.~\href{https://revistagi.geofisica.unam.mx/index.php/RGI/article/download/1638/1789/1961}{Link to download the document}} are as follows:
\begin{gather}
    \log\qty[N(M)] = 4.54 - 0.55 M.\label{eqn:numer-magnitude}
\end{gather}
Where $N(M)$ represents the number of earthquakes at the same magnitude, $M$, as defined on the Richter scale.

These models are essential for understanding the relationship between earthquake magnitude and the amount of energy released. 
For example, an increase of one unit on the Richter scale implies that approximately 32 times more energy is released. 
Therefore, a \num{7.0} magnitude earthquake releases energy equivalent to the explosion of \num{500 000} tons of TNT, whereas a \num{5.0} magnitude earthquake releases energy comparable to \num{500} tons of TNT.
The algebraic expression that relates the energy released by an earthquake with the following equation,
\begin{gather}
    \log\qty[E] = 4.8 + 1.5 M.\label{eqn:energy-magnitude}
\end{gather}
Where $M$ is the magnitude of the earthquake in the Richter scale and $E$ is the energy in joules (force times distance).
In this way, mathematics helps us not only to quantify natural phenomena, but also to prepare for their possible effects.

\section{Activity}

In this activity, we will use the given information to develop an understanding of how functions can be used to quantify natural phenomena.
To do this we are going to make algebraic manipulations to the equations~\eqref{eqn:numer-magnitude} and~\eqref{eqn:energy-magnitude} to find the following:
\begin{itemize}
    \item Find the algebraic expression for the difference of energy between different magnitudes.
\item Calculate the energy of an earthquake in terms of the frequency of earthquakes of the same magnitude in Mexico.
    \item Find a relation between the energy and the number of earthquakes with the same magnitude.
    \item Graph the functions and analize the order of magnitude of the energy.
\end{itemize}
In order to accomplish does objectives, fulfill the tasks below.

\subsection{Exploring with algebra}

First, we will find the algebraic expression for the difference in energy between earthquakes of different magnitudes.
Then we are going to express this difference in terms of energy and number of earthquakes with the same magnitude.
\begin{enumerate}
    \item Solve for $N(M)$ in equation~\eqref{eqn:numer-magnitude}
    \item Solve for $E$ in equation~\eqref{eqn:energy-magnitude}
    \item Compute the difference $E_2 - E_1$ using the previous result. You should arrive to \[E_2-E_1 = 10^{4.8}\qty[10^{1.5 M_2} - 10^{1.5 M_2}].\]
    \item Compute the difference between $\log\qty[E_2] - \log\qty[E_1]$. You should get \[\frac{E_2}{E_2} = 10^{1.5(M_2-M_1)}.\]
    \item Compute the difference $M_2-M_1$ using equation~\eqref{eqn:energy-magnitude}. You should get \[M_2-M_1 = \log\qty[\qty(\frac{E_2}{E_1})^{3/2}]\]
    \item Compute the difference $M_2-M_1$ using equation~\eqref{eqn:numer-magnitude}. You should get \[M_2-M_1 = \log\qty[\qty(\frac{N(M_2)}{N(M_1)})^{-1/0.55}].\]
    \item Using the two previous results, arrive at the following equation \[E_2 = \qty(\frac{N(M_2)}{N(M_1)})^{-2/(3\cdot 0.55)}E_1\]
    \item Considering that 1 Ton of TNT releases \num{4.184d9} joules, use equation~\eqref{eqn:energy-magnitude}, to find the magnitude of an earthquake that releases the same amount of energy.
\end{enumerate}

\subsection{Exploring with graphs}

Now we are going to explore the numeric values of the previous expressions.
\begin{enumerate}
    \item Using equation~\eqref{eqn:energy-magnitude}, compute the energy of the earthquakes shown in table~\ref{tab:sismos-mexico}.
    \item Divide the previous result by \num{4.184d9}, to express the energy in tons of TNT.
    \item Graph the numbers using Excel or by hand. Use logarithmic scale in the y axis.
    \item Compute the energy difference between the earthequake of march 28 and august 23. Express the energy in tons of TNT.
    \item Graph the equation~\eqref{eqn:energy-magnitude} for the earthquakes with magnitude $\{0,1,2,3,4,5,6,7,8,9,10\}$. Express the graph in terms of tons of TNT and use logarithmic scale in the $y$ axis.
    \item Identify the range and domain of equation~\eqref{eqn:energy-magnitude} using the previous graph.
\end{enumerate}

\subsection{Delivery instructions}

Upload a PDF document to Canvas containing the algebraic procedures.
Please enclose your final answer in a box.
Add the graphs and tables from the second part to the same document.
Although this activity is to be completed individually, you can work in groups.
Remember to include a cover page, and ensure your procedures are clear and orderly.

\subsection{Algebraic tools}

Finally, here are some algebraic properties that might be helpfull for the activity.

\begin{minipage}[c]{0.45\textwidth}
    \centering
    {\bfseries\large Laws of exponents}
    \begin{align*}
        a^r\cdot a^s = a^{r+s}, \quad & \qty(a^r)^s = a^{rs} \\
        \frac{a^r}{a^s} = a^{r-s}, \quad & \qty(ab)^r = a^r b^r \\
        \qty(\frac{a}{b})^r = \frac{a^r}{b^r},\quad & b\neq 0
    \end{align*}
\end{minipage}
\begin{minipage}[c]{0.45\textwidth}
    \centering
    {\bfseries\large Laws of Logarithms}
    \begin{align*}
        \log_a\qty(AB) &= \log_a\qty(A) + \log_a\qty(B) \\
        \log_a\qty(\frac{A}{B}) &= \log_a\qty(A) - \log_a\qty(B) \\
        \log_a\qty(A^C) &= C\log_a\qty(A)
    \end{align*}
\end{minipage}

\begin{minipage}[c]{\textwidth}
    \centering
    {\bfseries\large Inverse functions property}
    \begin{gather*}
        \qty(f\circ f^{-1})(x) = x = \qty(f^{-1}\circ f)(x)
        \quad\rightarrow\quad
        \log_a(a^x) = x = a^{\log_a(x)}
    \end{gather*}
\end{minipage}


\end{document}
